%----------
%	THESIS STRUCTURE
%----------	

La investigación se estructura de la siguiente manera en capítulos y secciones a lo largo de todo el documento:

\begin{table}[H]
    \ttabbox[\FBwidth]{
        \caption{Estructura de la Investigación}
        \label{table:thesis_structure} % Unique label for the table
    }{
        \begin{tabular}{p{3cm} p{11cm}}
            \hline
            \rowcolor{gray!30} % Adding a gray color to the header row
            \textbf{Capítulo / Anexo} & \textbf{Descripción} \\ \hline
            Capítulo 2: Estado del Arte & Revisión de la literatura sobre la representación mediática de la mujer, los sesgos de género y los métodos de detección basados en IA y procesamiento de lenguaje natural, así como los modelos de lenguaje, las agentes y las habilidades agénticas sobre los que se construye la solución, y la justificación de la propuesta. \\ \hline
            
            Capítulo 3: Metodología & Detalla el marco metodológico: el análisis experto con perspectiva de género, el corpus del proyecto IRIS\_IAMEDIA, la selección de las cinco variables de lenguaje y los modelos de lenguaje empleados. \\ \hline

            Capítulo 4: Implementación & Describe el desarrollo técnico de la solución: la arquitectura de agentes especializadas, el empaquetado de la metodología como Agent Skills, la recuperación de información (RAG) y el protocolo de evaluación, en línea con los objetivos de la Sección~\ref{sec:objectives_motivation}. \\ \hline

            Capítulo 5: Resultados y Discusión & Presenta los resultados de la evaluación variable a variable, la discusión de sus causas, repartidas entre el límite de la tarea y el de los modelos, y la aplicabilidad editorial mediante la integración en la herramienta de IRIS. \\ \hline
            
            Capítulo 6: Conclusiones y Líneas Futuras & Sintetiza los hallazgos principales de la investigación, reflexiona sobre el cumplimiento de los objetivos y propone futuras vías de desarrollo y aplicación. \\ \hline
            
            Anexo A: Presupuesto & Desglosa los costes humanos, tecnológicos y de licencias de \textit{software} necesarios para la viabilidad y ejecución de este proyecto. \\ \hline

    
            Anexo B: Código & Proporciona el acceso al repositorio digital que aloja el código fuente y la documentación técnica del sistema desarrollado. \\ \hline


            Anexo C: Cumplimiento del AI Act & Evalúa y justifica la exención del proyecto ante las normativas más restrictivas de la Ley de IA de la Unión Europea, por su carácter de investigación académica. \\ \hline
            Anexo D: Declaración de uso de IAG & Recoge la declaración de uso de Inteligencia Artificial Generativa cumplimentada conforme al modelo de la Universidad, con el uso técnico realizado y una reflexión sobre su utilidad. \\ \hline
        \end{tabular}
    }
\end{table}

\newpage

